\chapter{序論}
\label{chap:intro}

\section{背景}
\label{sec:intro-background}

拡散モデルは、データにノイズを徐々に加える過程を逆向きにたどることでデータを生成するモデルであり、画像生成で高い生成品質を示している~\cite{ho2020ddpm}。
Songら~\cite{song2021sde}は、この生成過程を確率微分方程式（SDE）と、それと同じ周辺分布をもつ常微分方程式（ODE）として定式化した。
Karrasら~\cite{karras2022edm}は、この枠組みのもとで学習とサンプリングの設計を整理し、CIFAR-10~\cite{krizhevsky2009cifar}などのデータセットで高い生成品質を報告している。
本論文では、Karrasらの学習方法とサンプリング方法をEDMと呼ぶ。

一方で、拡散モデルの学習には多くの計算を要する。
EDMのCIFAR-10の設定では、延べ2億枚の画像を用いてモデルを学習し、Karrasらはこの学習に8台のGPUを用いている。
利用できる計算資源が限られる場合、この学習量に達するまで学習を続けることは容易ではなく、学習を途中で打ち切ったモデル、すなわち学習不足のモデルを用いることになる。
また、学習に用いることのできるデータが少ない場合にも、生成品質の十分なモデルが得られるとは限らない。
このようなモデルの生成品質を、モデルを学習し直さずに改善できるかが、本研究の出発点となる問題である。

学習済みのモデルの生成品質は、サンプリング方法によっても変わる。
ODEに基づくサンプリングは少ないステップ数で画像を生成できる。
SDEに基づくサンプリングは、生成の途中でノイズを加えることで生成過程の誤差を抑える効果をもつが、どちらが優れるかは条件によって異なる~\cite{karras2022edm,cao2023odesde}。
Xuら~\cite{xu2023restart}のRestart Samplingは、ODEによる生成の途中でノイズを再び付与し、同じ区間の生成をやり直すことを繰り返す手法である。
この手法は、ODEの離散化誤差の小ささと、ノイズによる蓄積誤差の収縮を両立させることを狙っている。

Restart Samplingはモデルの再学習を必要としないため、学習不足のモデルの生成品質をサンプリング時に改善する手段になりうる。
しかし、Xuらの評価は公開されている学習済みモデルを対象としており、モデルの学習量によってRestart Samplingの効果や、区間をやり直す回数（以下、Restart回数$K$）の適切な値がどう変わるかは調べられていない。

\section{研究課題}
\label{sec:intro-rq}

本研究では、CIFAR-10のクラス条件付き生成を対象として、EDMの学習の途中で保存したモデルにRestart Samplingを適用し、次の三つの研究課題（RQ）を調べる。

\begin{description}
  \item[RQ1] 学習不足のモデルの生成品質を、Restartで改善できるか。
  \item[RQ2] 最適なRestart回数$K$は学習量にどう依存するか。
  \item[RQ3] 学習データが少ない場合も同じことが成り立つか。
\end{description}

ここで学習量とは、学習に用いた画像の延べ枚数を指す。
生成品質はFID~\cite{heusel2017fid}で評価する。
Restartを行わない場合（$K=0$）のFIDを基準FIDと呼び、Restartによって基準FIDから減少した量を改善幅と呼ぶ。

RQ1とRQ2では、CIFAR-10の全学習画像（50,000枚）で学習したモデルを対象とし、学習量と$K$を変えてFIDを比較する。
RQ3では、学習画像を30,000枚に減らして学習したモデルを対象として同じ比較を行う。
学習量が不足する場合と、学習データが不足する場合は、別の問題として区別して扱う。

\section{本研究の貢献}
\label{sec:intro-contributions}

本研究は新しいサンプリング手法を提案するものではなく、既存のRestart Samplingの効果が現れる条件を実験によって調べるものである。
貢献は次の三点である。

\begin{enumerate}
  \item \textbf{学習量とRestartの効果の関係の定量化。}
    学習量の異なる10個のモデルについて、Restart回数とFIDの関係を調べた。
    基準FIDが大きいモデル、すなわち学習量の少ないモデルほど、Restartによる改善幅は大きかった。
    この改善は、生成時の乱数を変えても再現した。
    十分に学習したモデルでは改善幅は小さく、生成時の乱数によるばらつきと同程度であった。
  \item \textbf{最適なRestart回数の学習量依存性。}
    FIDが最小となる$K$は、学習が進むにつれて小さくなる傾向を示した。
    メインの生成ステップ数を固定した条件に加え、モデルの評価回数（NFE）を固定した条件でも、学習量の少ないモデルでは大きな$K$が、十分に学習したモデルでは小さな$K$が最良となった。
  \item \textbf{学習不足とデータ不足での効果の比較。}
    学習画像を30,000枚に減らして学習したモデルについて同じ実験を行い、全学習画像で学習したモデルと比較する。
    \todo{結果は未定。30,000枚の連続学習（実験0-6）とそのRestart評価（実験8-15）の完了後に記述する。}
\end{enumerate}

\section{本論文の構成}
\label{sec:intro-organization}

% TODO(要修正): 各章を追加した時点で、章番号を \ref による参照に置き換える。
本論文の構成は次のとおりである。
第2章では、拡散モデル、EDM、Restart Sampling、評価指標FIDと、関連研究を説明する。
第3章では、データセットとモデル、学習条件、Restartの設定、評価手順を述べる。
第4章では、メインの生成ステップ数を固定した条件で、学習量ごとにRestart回数を変えた主実験の結果を示し、RQ1とRQ2を検証する。
第5章では、NFEを固定した条件で同じ比較を行う。
第6章では、学習データが少ない場合の結果を示し、RQ3を検証する。
\todo{第6章の内容（30,000枚での学習、全データでの延長学習との比較）は、実験0-6・0-7・8-15の完了後に確定する。}
第7章では、Restartの効果を左右する要因、$K$の選び方、本研究の限界を考察する。
第8章では、各研究課題への回答と今後の課題をまとめる。
% TODO(要修正): 付録A・Bの構成案は乱数再生のある実験（8-8の170k、8-1〜8-3、実験4）を含む。採否を決めてから次の文を確定する。
付録では、パラメータの選定、予備実験、周辺実験、再現条件を示す。
